(1) 对基模高斯解定义复曲率量
\[
\frac1{q(z)}=\frac1{R_c(z)}-i\frac{\lambda}{\pi w(z)^2}.
\]
由 \(q=z+iz_R\) 得 \(1/q=(z-iz_R)/(z^2+z_R^2)\)，比较实虚部得到
\[
R_c(z)=z\left(1+\frac{z_R^2}{z^2}\right),\qquad
w(z)^2=\frac{\lambda}{\pi}\frac{z^2+z_R^2}{z_R}
=w_0^2\left(1+\frac{z^2}{z_R^2}\right).
\]
在 z=0，波前曲率半径理解为无穷大，局部波前平面。

(2) 对 \(z\gg zR\)，\(w(z)\sim w_0z/z_R\)，故强度 \(1/e^{2}\) 半径定义下的远场发散半角
\[
\theta_{\rm div}=\frac{w_0}{z_R}=\frac{\lambda}{\pi w_0},
\quad w_0\theta_{\rm div}=\lambda/\pi.
\]
这个乘积是理想 TEM00 光束的衍射受限关系，不应与任意截断或多模激光的束质因子等同。

(3) 令 z=zR，有
\[
w(z_R)=\sqrt2\,w_0,\qquad R_c(z_R)=2z_R.
\]
因此 Rayleigh 长度是一侧从束腰增长到面积加倍、半径放大 sqrt2 的距离；在 \(z=-zR\) 同样成立。束腰前后虽然强度轮廓相同，但曲率 R 的符号相反。

(4) 因 \(w^2=w_0^2(1+z^2/z_R^2)\)，其唯一最小值为 \(w_{0}\)，发生在 z=0。若 \(w_{0}\to \infty \) 且 \(\lambda \) 固定，则 \(zR\propto w_{0}^{2}\to \infty \)、\(\theta div\to 0\)，对于任意有限 z，\(w(z)/w_0\to1\)，波前曲率趋无穷大，局部近似平面波。
